\begin{theorem}[Theorem 7]
Let $n\ge2$, $K\ge1$ and $0<\delta<1$, suppose the Sobolev--Poincar\'e
inequality on balls holds in dimension $n$ (\noderef{CitedSobolevPoincareBall}),
and assume, all at the exponent $p=n$, the cited regularity result (C1)
(\noderef{CitedRegularityEpsPos}), the cited solvability result (C2)
(\noderef{CitedDirichletW1p}) and the cited Dirichlet result (C3) with
H\"older boundary data (\noderef{CitedDirichletHolder}) in the universal forms
described below, together with the cited H\"older estimates (C4)
(\noderef{CitedHolderOscillationMinimum}) and the cited Campanato embedding (C5) at
the exponent $n/(n-1)$ on \emph{every} ball
(\noderef{CitedCampanatoEmbedding}).  Then there is $\beta>0$, depending only
on $n$, $K$ and $\delta$, with the following property.  Let
$\Omega\subseteq\R^n$ be a domain, let $F\colon\Omega\to\R^n$ be a
$K$-quasiregular mapping with component gradients $G$
(\noderef{IsQuasiregularOn}) and distortion matrix $\tau(x)$
(\noderef{IsDistortionMatrix}) and assume the cited properties (C7)
(\noderef{CitedQuasiregularBasics}) on \emph{every} open
$\Omega_0\subseteq\Omega$.  Then
\begin{itemize}
\item[(i)] if $\tau$ is a continuous function, then $F\in C^\eta_{\mathrm{loc}}(\Omega)$
  for each $0\le\eta<1$; and
\item[(ii)] if $\tau$ is H\"older continuous with exponent $\delta$, then
  $F\in C^{1+\beta}_{\mathrm{loc}}(\Omega)$.
\end{itemize}
Following the Note for the formalization, ``$\tau$ continuous'' (resp.\
``H\"older continuous with exponent $\delta$'') means that $\tau$ agrees almost
everywhere on $\Omega$ with a continuous (resp.\ $\delta$-H\"older continuous)
matrix-valued function, and the local conclusions are stated on every nonempty
bounded open $\Omega'$ with $\overline{\Omega'}\subseteq\Omega$ which is a
bounded Lipschitz domain (\noderef{IsBoundedLipschitzDomain}) and for which the
cited Campanato embedding (C5) holds on $\Omega'$, at the exponent $p=n$ and
also at the exponent $n/(n-1)$ (\noderef{CitedCampanatoEmbedding}).  \emph{The
Lipschitz hypothesis on $\Omega'$ is the one Theorem~Theorem~7 of the thesis inherits from
Theorems~Theorem~3 of the thesis and~Theorem~4 of the thesis, which state it at l.~1274; it is carried
as an explicit predicate rather than left in the phrase ``Lipschitz domain'',
and it is what those two theorems now require of the set passed in their
$\Omega'$ slot.}  Of the two (C5) instances the first is
the instance Theorem~Theorem~4 of the thesis consumes and the second the instance
Theorem~Theorem~3 of the thesis consumes, and both are carried per set because (C5) mentions
$\Omega'$ and so cannot be hoisted above it.  A \emph{third} instance of (C5),
at the exponent $n/(n-1)$ on every ball, is carried at the head instead: it
mentions no set of the conclusion, it is the instance
\noderef{MinimumGradientHolder} consumes internally and binds before that node's
own existential, and so it cannot be attached to $\Omega'$.  It is the Note's own
(C5) instantiated at a family each member of which satisfies its hypothesis,
since a ball is a bounded Lipschitz domain, and so is neither a stronger form nor
vacuous; and it carries no numerical content, mentioning only $n$.  The
conclusions are then, componentwise:
each $F^i$ agrees a.e.\ on $\Omega'$ with a function in
$C^\eta(\overline{\Omega'})$ (\noderef{MemHolderOn}), resp.\ each $F^i$ agrees
a.e.\ on $\Omega'$, together with its gradient $G^i$, with a function
differentiable at every point of $\Omega'$ whose gradient is continuous and
$\beta$-H\"older on $\overline{\Omega'}$ (\noderef{MemC1HolderOn} at the pair of
sets $(\Omega',\overline{\Omega'})$), i.e.\ in $C^{1+\beta}(\Omega')$ with the
gradient's H\"older seminorm taken on $\overline{\Omega'}$ (l.~278).

No nondegeneracy hypothesis on $F$ is imposed: Theorem~Theorem~7 of the thesis as printed
carries none, and the statement is true for constant maps as well.  For a
constant $F$ all component weak gradients vanish, so the formal Jacobian $J_F$
is identically $0$ and the second branch of \noderef{IsDistortionMatrix} gives
$\tau\equiv I$, for which both conclusions hold trivially.

\medskip
\noindent\textbf{Correction reported.}  Following the paper's Note for the
formalization, $\beta$ depends only on $n$, $K$ and $\delta$; the print has a
lower-case ``$k$'' in ``depending only on $k$ and $\delta$'', which is the
distortion constant $K$.

The hypothesis \noderef{CitedSobolevPoincareBall} is the cited result (C6) of the
paper's Note for the formalization, carried here as an explicit argument because
the proof uses it --- it enters through the constants of \noderef{MinimumMorreyRegularity} and \noderef{MinimumGradientHolder} --- which is what the
Note permits for the results the thesis quotes without proof.  Although it is
bound before the existential for $\beta$, it introduces no dependence of
$\beta$ on any quantity outside the list the paper names: (C6) is a statement
about $n$ alone, and $n$ is already bound first, so the hypothesis carries no
numerical content of its own.

\medskip
\noindent\textbf{The cited results of \S4 and the exponent $p=n$.}  The cited
regularity result (C1) (\noderef{CitedRegularityEpsPos}), the cited solvability
result (C2) (\noderef{CitedDirichletW1p}) and the cited Dirichlet result (C3)
with H\"older boundary data (\noderef{CitedDirichletHolder}) are assumed for
\emph{every} field $A'\colon\R^n\to\R^n$ and all parameters
$\varepsilon',\gamma_0',\gamma_1'$ with $\gamma_0'>0$ and $\gamma_1'>0$ --- and,
for (C1), for every open $\Omega_0\subseteq\R^n$ as well --- which are the forms
these three results have in the paper's Note for the formalization, where they
``may be taken as explicit hypotheses \ldots\ where they are used''.  They are
used here, and only here, through Theorems 3 and 4: the paper reaches
Theorem~Theorem~7 of the thesis by the sentence at l.~1626--1627, ``which is of the type
considered in \S4, with $a=0$, $\varepsilon=0$ and $\lambda$, $\Lambda$
depending only on $K$.  From Theorems 3 and 4 we deduce'', and both
\noderef{MinimumMorreyRegularity} and \noderef{MinimumGradientHolder} carry
(C1)--(C3) in exactly this universal form, bound before their own existentials.
By pm-0019 the Note's cited list is closed, so none of them is derivable and
each must be a hypothesis of the statement that uses it.

There is no free exponent $p$ in this theorem: the functional of \S5 is
$v\mapsto(\tau(x)\nabla v,\nabla v)^{n/2}$, so every one of (C1)--(C5) is
instantiated at $p=n$, which satisfies the standing requirement $1<p$ because
$n\ge2$.  (C4) mentions only $n$ and $p$ and is therefore bound at the head
alongside (C1)--(C3); (C5) is bound three times, twice with $\Omega'$ --- once in
each of the two conclusions --- and once at the head in the on-balls form, which
mentions no set of the conclusion.  None of the five introduces a dependence of $\beta$ on
any quantity outside the list the paper names at l.~1635 (``$\beta>0$ depending
only on $k$ and $\delta$'', with $k$ the distortion constant $K$): each mentions
only $n$ and the exponent $n$, both already bound, and each is universally
quantified over its remaining parameters, so none carries numerical content of
its own --- the same reading the Note's own wording licenses for (C6) above.

\medskip
\noindent\textbf{(C7) in the universal form.}  The cited properties (C7)
(\noderef{CitedQuasiregularBasics}) are assumed for \emph{every} open
$\Omega_0\subseteq\Omega$, not only for $\Omega$ itself.  That is the form (C7)
has in the paper's Note, which states it for ``a $K$-quasiregular mapping $F$''
with no distinguished domain and whose local-minimum clause already quantifies
over the subsets $E$: the restriction of a $K$-quasiregular mapping to an open
subset is again $K$-quasiregular, so the family is the cited result applied to
restrictions, which is neither a stronger form nor vacuous --- each instance is
moreover \emph{conditional} on $F$ being $K$-quasiregular on $\Omega_0$ with
distortion matrix $\tau$ there (\noderef{IsQuasiregularOn},
\noderef{IsDistortionMatrix}), so no instance asserts anything about a set on
which the hypotheses fail.  It is needed because the proof consumes the
local-minimum clause on an interposed bounded open set $\Omega''$ with
$\overline{\Omega'}\subseteq\Omega''$ and
$\overline{\Omega''}\subseteq\Omega$, which is where the boundedness that
Theorems~Theorem~3 of the thesis and~Theorem~4 of the thesis require of their domain comes from, and
$\Omega''$ is not in general a ball.  Being bound inside the quantifiers over
$\Omega,F,G,\tau$ and after the existential for $\beta$, it carries no
numerical content and no dependence of $\beta$ on anything.

\textbf{On non-vacuity of the quasiregularity hypotheses.}  The statement is
conditioned on \noderef{IsQuasiregularOn} and \noderef{IsDistortionMatrix}, so it
would be vacuously true if no triple $(F,G,\tau)$ satisfied the two together,
which the Note for the formalization forbids at l.~258.  It does not:
\noderef{ModelQuasiregularIdentity} exhibits the identity map of $\R^n$, with
$G^i(x)=e_i$ and $\tau(x)=I$, as a $1$-quasiregular mapping with distortion
matrix $\tau$ on every open set --- the paper's own trivial instance of
(5.10) of the thesis, which the identity satisfies with equality at the least
admissible constant $K=1$.  The citation is for non-vacuity only; no step of the
proof below uses the identity map.
\end{theorem}
\begin{proof}
Throughout $p=n$ and $\varepsilon=0$, which is the instantiation the paper names
at l.~1626--1627 (``which is of the type considered in \S4, with $a=0$,
$\varepsilon=0$''), and $1<p$ holds because $n\ge2$.  For a matrix $S$ and
vectors $v,w\in\R^n$ write $(Sv,w)=\sum_a\sum_bS_{ab}v_aw_b$.  We write $\tau_M$
for the Morrey exponent that \noderef{MinimumMorreyRegularity} calls $\tau$, to
avoid the clash with the distortion matrix.

\medskip
\noindent\emph{Part 0: the constants, fixed before $\Omega,F,G,\tau$.}

\textbf{(0a)} Apply \noderef{TauPropertiesQuasiregular} with $n$ and $K$; its
hypotheses are $2\le n$ and $1\le K$, both data of the present theorem.  It
produces $\lambda,\Lambda$ with $0<\lambda\le\Lambda$, depending only on $n$ and
$K$ because they are bound immediately after $n,K$ and before all remaining
data.  Fix them.

\textbf{(0b)} Put
\[
\varepsilon=0,\qquad a=1,\qquad
\nu_0=\frac nK,\qquad \nu_1=2n(n-1)K^{\,n-1},\qquad \tau_M=\frac n2 .
\]
All four of $0\le\varepsilon\le1$, $0<a$, $0<\nu_0$ and $0<\tau_M<n$ hold: the
first two by inspection; $\nu_0>0$ because $n\ge2>0$ and $K\ge1>0$; and
$0<n/2<n$ because $n\ge2$.  For $\nu_0<\nu_1$, since $K>0$ the inequality
$n/K<2n(n-1)K^{n-1}$ is equivalent to $1<2(n-1)K^{\,n}$, and $n\ge2$ gives
$n-1\ge1$ while $K\ge1$ gives $K^{\,n}\ge1$, so $2(n-1)K^{\,n}\ge2>1$.

\emph{Why $\nu_1$ is twice the derived constant.}  The (F2) computation inside
\noderef{QuasiregularLocalVariationalData} produces the upper bound with the
constant $n(n-1)K^{\,n-1}$, which is why that node asks only for
$n(n-1)K^{\,n-1}\le\nu_1$.  That
constant is \emph{not} usable as $\nu_1$, because
\noderef{MinimumMorreyRegularity} and \noderef{MinimumGradientHolder} both
require $\nu_0<\nu_1$ \emph{strictly}, and at $n=2$, $K=1$ one has
$\nu_0=n/K=2$ and $n(n-1)K^{n-1}=2$, so the two would be equal and both nodes
would be inapplicable --- and $n=2$, $K=1$ is a legitimate instance of
Theorem~Theorem~7 of the thesis.  Doubling is free: (F2)'s upper half is an inequality, and
$2n(n-1)K^{n-1}\ge n(n-1)K^{n-1}$, so the derived bound implies the bound with
the doubled constant.  Both constants depend only on $n$ and $K$, as the theorem
requires.

\textbf{(0c)} Apply \noderef{MinimumGradientHolder} with $n$, $p=n$, the
$\lambda,\Lambda$ of (0a), $a=1$, $\nu_0,\nu_1$ of (0b) and the $\delta$ of the
present theorem.  Its head hypotheses are discharged thus: $2\le n$, $1<p$ (as
$p=n\ge2$), $0<\lambda$ and $\lambda\le\Lambda$ by (0a), $0<a$, $0<\nu_0$ and
$\nu_0<\nu_1$ by (0b), $0<\delta$ by hypothesis; the Sobolev--Poincar\'e
inequality \noderef{CitedSobolevPoincareBall} is a hypothesis of the present
theorem and is passed on unchanged; the three universal clauses (C1), (C2), (C3)
(\noderef{CitedRegularityEpsPos}, \noderef{CitedDirichletW1p},
\noderef{CitedDirichletHolder}) are hypotheses of the present theorem in exactly
the shape that node asks for, namely for every field $A'$, all
$\varepsilon',\gamma_0',\gamma_1'$ with $\gamma_0'>0$ and $\gamma_1'>0$ and, for
(C1), every open $\Omega_0$, all at the exponent $p=n$; and the on-balls
Campanato hypothesis --- (C5) at the exponent $n/(n-1)$ on $B_\rho(y)$ for every
$y$ and every $\rho>0$ (\noderef{CitedCampanatoEmbedding}) --- is the head
hypothesis of the present theorem, carried for this purpose.  The node produces
\[
\lambda_0\quad\text{with}\quad 0<\lambda_0\le p=n .
\]
Put $\beta=\lambda_0/n$.  Then $\beta>0$.  And $\beta$ depends only on $n$, $K$
and $\delta$: in \noderef{MinimumGradientHolder} the constant $\lambda_0$ is
bound immediately after $n,p,\lambda,\Lambda,a,\nu_0,\nu_1,\delta$ and the cited
hypotheses, and \emph{before} $c_\sigma,d,M,W_n,\sigma$ and before
$\varepsilon,\Omega,\Omega',f^0,u,g$; each of the quantities it may depend on is
a function of $n,K,\delta$ alone ($p=n$; $\lambda,\Lambda$ from (0a) at $(n,K)$;
$a=1$; $\nu_0,\nu_1$ from (0b); $\delta$ itself), and each cited hypothesis
mentions only $n$ and the exponent $n$ and is universally quantified over its
remaining parameters, so none carries numerical content.  This is the quantifier
order pm-0007 requires, and it is why $\beta$ may be --- and is --- fixed here,
before $\Omega,F,G,\tau$ are introduced.

Now let $\Omega,F,G,\tau$ be given with $\Omega$ open and connected, $F$
$K$-quasiregular on $\Omega$ with component gradients $G$
(\noderef{IsQuasiregularOn}), $\tau$ its distortion matrix
(\noderef{IsDistortionMatrix}), and (C7) (\noderef{CitedQuasiregularBasics})
available on every open $\Omega_0\subseteq\Omega$.  The argument splits on
whether $F$ is \emph{almost everywhere} constant on $\Omega$; no nondegeneracy
hypothesis is needed, matching the printed statement of Theorem~Theorem~7 of the thesis,
which carries none.  \emph{The split is taken almost-everywhere-invariantly
because that is the form in which both \textup{(C7)} and
\noderef{TauPropertiesQuasiregular} state their non-constancy condition}: $F$ is
a $W^{1,n}_{\mathrm{loc}}$ map, determined only up to a null set, so a pointwise
non-constancy condition would be a condition on the representative rather than
on the map, and the two branches below must between them exhaust the genuine
alternatives.  They do: either some $c$ has $F=c$ almost everywhere on $\Omega$,
or none does.

\medskip
\noindent\emph{Part 1: the case $F$ almost everywhere constant on $\Omega$.}

Suppose there is $c\in\R^n$ with $F=c$ almost everywhere on $\Omega$.  Fix $i$.
Apply \noderef{QuasiregularConstantCase} with the parameters $n$ and $K$ --- its
only guard is $2\le n$, a hypothesis of the present theorem --- on the open set
$\Omega$, with the present $F$ and $G$ and with this $c$; its remaining
hypotheses are the openness of $\Omega$ and \noderef{IsQuasiregularOn}, both
hypotheses of the present theorem, and $F=c$ almost everywhere on $\Omega$, which
is the branch hypothesis just assumed.  It gives, for the present $i$,
\[
G^i=0\quad\text{a.e.\ on }\Omega .
\tag{1.1}
\]
\emph{That node is where the almost-everywhere uniqueness of a weak gradient is
used}, and it is the only place in this proof that needs it; it is also the only
thing the constant branch needs beyond the branch hypothesis itself.

Both conclusions now follow with no cited input, and in particular without using
either $\tau$-antecedent.  Let $\Omega'$ be as in the conclusion, so
$\Omega'\subseteq\overline{\Omega'}\subseteq\Omega$.

For (i), let $0\le\eta<1$ and put $w(x)=c_i$ for all $x$.  Then
\[
F^i=c_i=w\qquad\text{almost everywhere on }\Omega' ,
\tag{1.2}
\]
which is the branch hypothesis restricted: it gives $F=c$ almost everywhere on
$\Omega$, hence $F^i=c_i$ almost everywhere on $\Omega$ componentwise, and
$\Omega'\subseteq\Omega$ as recorded above, so an almost-everywhere statement on
$\Omega$ holds almost everywhere on $\Omega'$ --- the exceptional set only
shrinks.  \emph{No pointwise identity is available or needed.}  $F$ is carried by
\noderef{IsQuasiregularOn} as a $W^{1,n}_{\mathrm{loc}}$ map, determined up to a
null set, and the branch hypothesis is itself an almost-everywhere statement, so
$F(x)^i=c_i$ can fail at points of $\Omega'$; what clause (i) of the conclusion
asks for is exactly the almost-everywhere identity $(1.2)$.  And
$w\in C^\eta(\overline{\Omega'})$ (\noderef{MemHolderOn}): $w$ is continuous, and
the H\"older constant $C=0$ works, since $|w(x)-w(y)|=0\le0\cdot\mathrm{dist}
(x,y)^\eta$ for all $x,y$.

For (ii), put $w(x)=c_i$ and $g_w(x)=0$ for all $x$.  Then $F^i=w$ almost
everywhere on $\Omega'$, which is $(1.2)$ again, and $G^i=0=g_w$ almost
everywhere on $\Omega'$ by (1.1), since $\Omega'\subseteq\Omega$ and the same
restriction of an almost-everywhere statement applies.  These are the two
almost-everywhere identities clause (ii) of the conclusion asks for; as in (i),
no pointwise identity for $F^i$ is available.  And
$\mathrm{MemC1HolderOn}\;n\;\beta\;\Omega'\;\overline{\Omega'}\;w\;g_w$
(\noderef{MemC1HolderOn}) holds, the three clauses in turn: $w$ is
differentiable at every point of $\R^n$, in particular at every point of
$\Omega'$, with vanishing differential, so each directional derivative equals
the corresponding component of $g_w=0$; $g_w$ is continuous, hence continuous on
$\overline{\Omega'}$; and the constant $C=0$ works for the H\"older seminorm of
$g_w$ on $\overline{\Omega'}$, because $\|g_w(x)-g_w(y)\|=0$.  \emph{The
constant branch is insensitive to which pair of sets is taken}, since all three
clauses hold for $w=c_i$, $g_w=0$ on every subset of $\R^n$.  This
holds for the $\beta>0$ fixed in (0c), which depends only on $n$, $K$ and
$\delta$.

\emph{Neither branch is vacuous on constant maps.}  Here every component weak
gradient vanishes a.e.\ by (1.1), so the formal differential $DF=(G^i_j)_{ij}$
is the zero matrix a.e.\ and $J_F=\det DF=0$ a.e.  At every $x\in\Omega$ where
$DF(x)=0$ we have $\neg(0<J_F(x))$, so the second branch of
\noderef{IsDistortionMatrix} applies and $\tau(x)=I$.  Thus $\tau=I$ a.e.\ on
$\Omega$, and $I$ is continuous and $\delta$-H\"older for every $\delta$ (with
constant $0$), so both antecedents are satisfiable; $I$ also satisfies the
corrected two-sided bound $K^{-2/n}|\xi|^2\le|\xi|^2\le K^{2(n-1)/n}|\xi|^2$
whenever $K\ge1$.

\medskip
\noindent\emph{Part 2: the case $F$ not almost everywhere constant on $\Omega$
--- the common setup.}

Assume now that there is no $c$ with $F=c$ almost everywhere on $\Omega$, which
is exactly the nonconstancy hypothesis of
\noderef{TauPropertiesQuasiregular}, and exactly the antecedent of the Jacobian
clause of (C7) (\noderef{CitedQuasiregularBasics}); both state it
almost-everywhere-invariantly, because $F$ is determined only up to a null set.

\textbf{(2a)} Apply the (C7) hypothesis at $\Omega_0=\Omega$, legitimate since
$\Omega$ is open and $\Omega\subseteq\Omega$, to get
$\mathrm{CitedQuasiregularBasics}\ n\ K\ \Omega\ F\ G\ \tau$
(\noderef{CitedQuasiregularBasics}).  Feeding it the present theorem's
\noderef{IsQuasiregularOn} and \noderef{IsDistortionMatrix} hypotheses yields its
four conclusions on $\Omega$.  We use clause (1), which supplies a field
$\widetilde F$ with $\widetilde F$ continuous on $\Omega$ and $F=\widetilde F$
almost everywhere on $\Omega$; and clause (2), which under the nonconstancy just
assumed gives $0<J_F$ almost everywhere on $\Omega$.  Clause (3), the
weak-solution property, is not used.

\emph{Clause (1) does not say that $F$ is continuous, and this is used below at
(3e).}  $F$ is carried by \noderef{IsQuasiregularOn} only as a
$W^{1,n}_{\mathrm{loc}}$ map, determined up to a null set, so continuity is a
property of the representative $\widetilde F$ and not of $F$: redefining $F$ at
one point of $\Omega$ leaves every hypothesis of the present theorem intact and
destroys continuity.  Fix such an $\widetilde F$ once and for all.

\textbf{(2b)} Apply \noderef{TauPropertiesQuasiregular} --- the instance of (0a),
at the same $n,K$ --- with the data $\Omega,F,G,\tau$.  Its hypotheses are: the
openness of $\Omega$; \noderef{IsQuasiregularOn} and
\noderef{IsDistortionMatrix}, both hypotheses of the present theorem;
$\mathrm{CitedQuasiregularBasics}\ n\ K\ \Omega\ F\ G\ \tau$, which is (2a); and
the nonconstancy of $F$ on $\Omega$, which is the case assumption.  It gives:
for almost every $x\in\Omega$, $\tau(x)$ is symmetric, $\det\tau(x)=1$, and
\[
K^{-\frac2n}|\xi|^2\le(\tau(x)\xi,\xi)\le K^{\frac{2(n-1)}n}|\xi|^2
\quad\text{for all }\xi\in\R^n,
\tag{2.1}
\]
\[
\lambda|h|^{\,n}\le(\tau(x)h,h)^{\frac n2}\le\Lambda|h|^{\,n}
\quad\text{for all }h\in\R^n .
\tag{2.2}
\]

\emph{The upper bound in \textup{(2.1)} is $K^{2(n-1)/n}|\xi|^2$ and not
$|\xi|^2$.}  The thesis prints $(\tau(x)\xi,\xi)\le|\xi|^2$ at l.~1608--1610, and
the manuscript's own \verb|% CHECK| at l.~1603--1605 records that this is false
in general: $\det\tau=1$ forces the largest eigenvalue of $\tau$ to be at least
$1$, with equality only for $\tau=I$ --- the counterexample given there is $n=2$,
$DF=\mathrm{diag}(a,b)$ with $a>b$, for which $\tau=\mathrm{diag}(b/a,a/b)$ has
largest eigenvalue $a/b>1$ --- and the same note prescribes the two-sided bound
of (2.1), adding that any $K$-dependent bounds suffice.  Accordingly
\noderef{TauPropertiesQuasiregular} states (2.1) and nothing stronger, and (2.1)
is what this proof uses; this is why $\nu_1$ below carries a power of $K$, in
agreement with l.~1617--1618's ``constants depending exclusively on $K$''.

\medskip
\noindent\emph{Part 3: the local data attached to a given $\Omega'$, from
\noderef{QuasiregularLocalVariationalData}.}

Both conclusions have the same shape: a set $\Omega'$ that is open, nonempty and
bounded with $\overline{\Omega'}\subseteq\Omega$, carrying (C5) at the exponents
$n$ and $n/(n-1)$ (\noderef{CitedCampanatoEmbedding}), and a representative
$\tau'$ of $\tau$.  Fix such an $\Omega'$ and such a $\tau'$, so that
\[
\tau=\tau'\ \text{a.e.\ on }\Omega,
\tag{3.1}
\]
and assume for the whole of Part 3 only this about $\tau'$: \emph{every entry
$x\mapsto\tau'(x)_{jk}$ is continuous on $\Omega$.}  In part (i) that is the
antecedent verbatim.  In part (ii) the antecedent gives $C\ge0$ with
$|\tau'(x)_{jk}-\tau'(y)_{jk}|\le C\,\mathrm{dist}(x,y)^\delta$ for all
$x,y\in\Omega$, and since $\delta>0$ this forces continuity on $\Omega$: given
$x\in\Omega$ and $\epsilon>0$, take the radius $(\epsilon/C)^{1/\delta}$ when
$C>0$ and any radius when $C=0$, and for $y\in\Omega$ within it
$|\tau'(y)_{jk}-\tau'(x)_{jk}|\le C\,\mathrm{dist}(y,x)^\delta<\epsilon$, the
strict inequality because $r\mapsto r^{\delta}$ is strictly increasing on
$[0,\infty)$ for $\delta>0$.  So Part 3 applies in both cases.  Fix also
$i\in\{1,\dots,n\}$ and write $u=F^i$, $g=G^i$.

Apply \noderef{QuasiregularLocalVariationalData} with the parameters $n$, $K$,
the $\lambda,\Lambda$ of (0a) and the $\nu_0,\nu_1$ of (0b).  Its head guards
hold: $2\le n$ and $1\le K$ are data of the present theorem; $0<\lambda$ and
$\lambda\le\Lambda$ are (0a); $0<\nu_0$ is (0b), and $\nu_0\le n/K$ holds with
equality by the choice $\nu_0=n/K$ there; and $n(n-1)K^{\,n-1}\le\nu_1$ is
exactly the doubling of (0b), since $\nu_1=2n(n-1)K^{\,n-1}$ and
$n(n-1)K^{\,n-1}\ge0$.  Its data hypotheses are, in order: the openness of
$\Omega$ and of $\Omega'$, the nonemptiness and the boundedness of $\Omega'$ and
$\overline{\Omega'}\subseteq\Omega$, all data of the conclusion being proved;
\noderef{IsQuasiregularOn} and \noderef{IsDistortionMatrix}, hypotheses of the
present theorem; the (C7) hypothesis in its universal form
(\noderef{CitedQuasiregularBasics}), a hypothesis of the present theorem in
exactly that shape; the nonconstancy of $F$ on $\Omega$, which is the case
assumption of Part 2; the almost-everywhere equality $(3.1)$ and the continuity
of the entries of $\tau'$, both just fixed; and the almost-everywhere two-sided
bounds, which are $(2.1)$ and $(2.2)$ of (2b).

It produces $\Omega''$, $f^0$, $\sigma$, $M$, $W_n$ and $d$ with
\[
\overline{\Omega'}\subseteq\Omega'',\qquad
\overline{\Omega''}\subseteq\Omega,\qquad
d\le\mathrm{dist}(x,y)\quad(x\in\Omega',\ y\notin\Omega''),
\tag{3.3}
\]
$\Omega''$ open, and $0<d$, $0\le M$, $0\le W_n$; with
\[
\mathrm{VariationalStructureF123}\ n\ n\ 0\ \lambda\ \Lambda\ 1\ \nu_0\ \nu_1\
\Omega''\ f^0\ \sigma
\tag{3i}
\]
(\noderef{VariationalStructureF123}); with
\[
u\in W^{1,n}(\Omega'')\ \text{with weak gradient}\ g
\tag{3.9}
\]
(\noderef{MemW1pOn}); with $|u|\le M$ almost everywhere on $\Omega''$ and
$\int_{\Omega''}|u|^n\,dx+\int_{\Omega''}\|g\|^n\,dx\le W_n^{\,n}$, which we
refer to below as (3e); with
\[
\mathrm{IsLocalMinimumOn}\ n\ n\ \Omega''\ f^0\ u\ g
\tag{3.8}
\]
(\noderef{IsLocalMinimumOn}); and with the conditional H\"older rate: for every
$\delta$ with $0<\delta<1$ and every $C\ge0$ such that
$|\tau'(x)_{jk}-\tau'(y)_{jk}|\le C\,\mathrm{dist}(x,y)^\delta$ for all
$x,y\in\Omega$ and all $j,k$, there is $c_\sigma>0$ with
$\sigma(t)\le c_\sigma t^\delta$ for every $t\ge0$.  \emph{That last clause is
what Part 5 consumes and Part 4 does not}, which is why it is a conditional
inside the node's conclusion rather than a hypothesis of it.

\emph{The node's own $\Omega$ slot is the present $\Omega$ and its $\Omega'$ slot
the present $\Omega'$; the interposed $\Omega''$ is produced, not supplied.}  In
Parts 4 and 5 it is $\Omega''$ that occupies the ``$\Omega$'' slot of the two
\S4 theorems and the present $\Omega'$ that occupies their ``$\Omega'$'' slot, so
the Lipschitz and (C5) hypotheses those theorems place on that slot are the data
of the conclusion being proved and nothing is asked of $\Omega''$ beyond
$(3.3)$.

\medskip
\noindent\emph{Part 4: conclusion \textup{(i)}.}

Assume the first antecedent: $\tau$ agrees a.e.\ on $\Omega$ with a $\tau'$ whose
entries are continuous on $\Omega$.  Let $\Omega'$ and $\eta$ be given with
$\Omega'$ open, nonempty and bounded, $\overline{\Omega'}\subseteq\Omega$, (C5)
on $\Omega'$ at the exponents $n$ and $n/(n-1)$, and $0\le\eta<1$; fix $i$.  If
$F$ is almost everywhere constant on $\Omega$ we are done by Part 1, so assume
not and let Part 2 apply; the two branches are exhaustive because the branch
condition is a disjunction of a statement and its negation.  \emph{The order
matters: Part 3 consumes $(2.1)$, $(2.2)$ and (2a), which exist only under the
nonconstancy assumption, so the branch must be taken first.}  Now run Part 3,
obtaining $\Omega''$, $f^0$, $M$, $W_n$, $d$, $\sigma$ and the conclusions
(3.8), (3.9) and (3i).

Apply \noderef{MinimumMorreyRegularity} with $n$, $p=n$, $\tau_M=n/2$, the
$\lambda,\Lambda$ of (0a), $a=1$, $\nu_0,\nu_1$ of (0b), and $d$, $M$, $W_n$,
$\sigma$ of Part 3.  Its hypotheses, in order:
\begin{itemize}
\item $2\le n$, $1<p$ (as $p=n\ge2$), $0<\tau_M$ and $\tau_M<n$ (as
  $\tau_M=n/2$ and $n\ge2$), $0<\lambda$, $\lambda\le\Lambda$, $0<a$, $0<\nu_0$,
  $\nu_0<\nu_1$: by (0a) and (0b).  \emph{The exponent $\tau_M$ must be
  positive}, which is why $n/2$ and not $0$ is passed: the node binds
  $0<\tau_M<n$, matching the manuscript's \verb|% CHECK| at l.~1277 and l.~1554,
  which records that $\tau_M=0$ must be excluded because at $\tau_M=0$ clause (i)
  would assert $|\nabla u|\in L^{\infty}$ and the proof needs $\tau_M>0$.  The
  clause of the conclusion used below does not mention $\tau_M$, so any
  admissible value serves.
\item $0<d$, $0\le M$, $0\le W_n$: by (3e).
\item \noderef{CitedSobolevPoincareBall}, and the three universal clauses (C1),
  (C2), (C3) at the exponent $p=n$ (\noderef{CitedRegularityEpsPos},
  \noderef{CitedDirichletW1p}, \noderef{CitedDirichletHolder}): hypotheses of the
  present theorem, in exactly that shape, passed on unchanged.
\item $0\le\varepsilon$ and $\varepsilon\le1$: at $\varepsilon=0$.
\item $\Omega''$ and $\Omega'$ open, $\Omega'$ nonempty and bounded,
  $\overline{\Omega'}\subseteq\Omega''$: by (3.3) and the data of the conclusion.
  Note the node's own ``$\Omega$'' is instantiated as $\Omega''$, not as
  $\Omega$.
\item $\forall x\in\Omega'$, $\forall y\notin\Omega''$,
  $d\le\mathrm{dist}(x,y)$: by (3e).
\item $\Omega'$ a bounded Lipschitz domain
  (\noderef{IsBoundedLipschitzDomain}): a datum of the conclusion.  \emph{The
  node's $\Omega'$ slot is the present $\Omega'$, not $\Omega''$}, so the
  hypothesis Theorem~3 places on that slot is supplied by the conclusion's own
  datum and nothing is asked of the interposed set.
\item (C5) at the exponent $n/(n-1)$ on $\Omega'$
  (\noderef{CitedCampanatoEmbedding}): a datum of the conclusion.
\item (C4) at the exponent $p=n$ (\noderef{CitedHolderOscillationMinimum}): a hypothesis
  of the present theorem.
\item $\mathrm{VariationalStructureF123}$ at
  $n,n,0,\lambda,\Lambda,1,\nu_0,\nu_1,\Omega'',f^0,\sigma$: by (3i).
\item $u\in W^{1,n}(\Omega'')$ with weak gradient $g$: by (3.9).
\item $|u|\le M$ a.e.\ on $\Omega''$, and
  $\int_{\Omega''}|u|^n+\int_{\Omega''}\|g\|^n\le W_n^{\,n}$: by (3e).
\item $\mathrm{IsLocalMinimumOn}\ n\ n\ \Omega''\ f^0\ u\ g$: by (3.8).
\end{itemize}
Every hypothesis is discharged, so the conclusion holds; we use its third clause,
which for every $\eta$ with $0\le\eta<1$ gives a $w$ with $u=w$ a.e.\ on
$\Omega'$ and $w\in C^\eta(\overline{\Omega'})$ (\noderef{MemHolderOn}).  As
$u=F^i$ and $i$ was arbitrary, this is conclusion (i).

\medskip
\noindent\emph{Part 5: conclusion \textup{(ii)}.}

Assume the second antecedent: there are $\tau'$ and $C\ge0$ with
$\tau=\tau'$ a.e.\ on $\Omega$ and
$|\tau'(x)_{ij}-\tau'(y)_{ij}|\le C\,\mathrm{dist}(x,y)^\delta$ for all
$x,y\in\Omega$ and all $i,j$.  Let $\Omega'$ be given as in the conclusion; fix
$i$.  Again, if $F$ is almost everywhere constant on $\Omega$ we are done by
Part 1, so assume not, and Part 2 applies; as in Part 4 the branch is taken
before Part 3, which consumes data available only under the nonconstancy
assumption.  Since $\delta>0$, the H\"older bound makes every entry of $\tau'$
continuous on $\Omega$, so Part 3 applies verbatim and yields
$\Omega''$, $f^0$, $M$, $W_n$, $d$, $\sigma$, (3.8), (3.9) and (3i).

In addition, $\sigma$ satisfies the H\"older condition (F4).  This is the
conditional clause of \noderef{QuasiregularLocalVariationalData} recorded at the
end of Part 3, read at the present $\delta$ and at the $C$ of the antecedent: its
three guards are $0<\delta$ and $\delta<1$, hypotheses of the present theorem,
and $0\le C$, part of the antecedent, and its hypothesis is the H\"older bound
$|\tau'(x)_{ij}-\tau'(y)_{ij}|\le C\,\mathrm{dist}(x,y)^\delta$ for all
$x,y\in\Omega$ and all $i,j$, which is the antecedent verbatim.  It yields
$c_\sigma>0$ with
\[
\sigma(t)\le c_\sigma t^\delta\qquad\text{for all }t\ge0 .
\]
\emph{The strictness $0<c_\sigma$ matters and is delivered by that node}, because
\noderef{MinimumGradientHolder} requires it while the antecedent here allows
$C=0$.

Apply \noderef{MinimumGradientHolder} --- the instance of (0c), at the same
$n,p=n,\lambda,\Lambda,a=1,\nu_0,\nu_1,\delta$, so that the $\lambda_0$ it
produces is the one fixed there and $\beta=\lambda_0/n$ --- with
$c_\sigma$, $d$, $M$, $W_n$ and $\sigma$.  The head hypotheses were discharged in
(0c).  The remaining ones, in order:
\begin{itemize}
\item $0<c_\sigma$, $0<d$, $0\le M$, $0\le W_n$: just above, and (3e).
\item $0\le\varepsilon$, $\varepsilon\le1$: at $\varepsilon=0$.
\item $\Omega''$ and $\Omega'$ open, $\Omega'$ nonempty and bounded,
  $\overline{\Omega'}\subseteq\Omega''$: by (3.3) and the data of the
  conclusion; again the node's ``$\Omega$'' is $\Omega''$.
\item $\forall x\in\Omega'$, $\forall y\notin\Omega''$,
  $d\le\mathrm{dist}(x,y)$: by (3e).
\item $\Omega'$ a bounded Lipschitz domain
  (\noderef{IsBoundedLipschitzDomain}): a datum of the conclusion, and again it
  is the present $\Omega'$ and not $\Omega''$ that occupies the node's $\Omega'$
  slot.
\item (C5) on $\Omega'$ at the exponent $p=n$ \emph{and} at the exponent
  $n/(n-1)$ (\noderef{CitedCampanatoEmbedding}): both are data of the conclusion.
\item (C4) at $p=n$ (\noderef{CitedHolderOscillationMinimum}): a hypothesis of the
  present theorem.
\item $\mathrm{VariationalStructureF123}$ at
  $n,n,0,\lambda,\Lambda,1,\nu_0,\nu_1,\Omega'',f^0,\sigma$: by (3i).
\item $\sigma(t)\le c_\sigma t^\delta$ for all $t\ge0$: just above.
\item $u\in W^{1,n}(\Omega'')$ with weak gradient $g$: by (3.9).
\item $|u|\le M$ a.e.\ on $\Omega''$ and the $L^n$ bound by $W_n^{\,n}$: by (3e).
\item $\mathrm{IsLocalMinimumOn}\ n\ n\ \Omega''\ f^0\ u\ g$: by (3.8).
\end{itemize}
Every hypothesis is discharged.  The second clause of the conclusion gives $w$
and $g_w$ with $u=w$ a.e.\ on $\Omega'$, $g=g_w$ a.e.\ on $\Omega'$ and
$\mathrm{MemC1HolderOn}\;n\;(\lambda_0/p)\;\Omega'\;\overline{\Omega'}\;w\;g_w$
(\noderef{MemC1HolderOn}) --- $w$ differentiable at every point of $\Omega'$
with $\nabla w=g_w$ there, and $g_w$ continuous and
$(\lambda_0/p)$-H\"older on $\overline{\Omega'}$.  That is the pair of sets the
present conclusion (ii) asks for, so it is handed on unchanged.  Since $p=n$, the exponent is $\lambda_0/n=\beta$, the
constant fixed in (0c), which depends only on $n$, $K$ and $\delta$.  As $u=F^i$,
$g=G^i$ and $i$ was arbitrary, this is conclusion (ii).
\end{proof}
